\chapter[Conclusion]{Conclusion}

The overall conclusion of the proposed project is presented here. The work carried out is summarised, the major achievements of the system are highlighted, and the extent to which the project objectives were achieved is discussed. Learning outcomes from the project are also presented, along with identified limitations and a roadmap for future enhancements.

\section[Conclusion]{Conclusion}
This project addressed the critical problem of supply chain volatility and the lack of automated resilience in small and medium enterprise (SME) logistics networks through the development of an Agentic Supply-Chain Digital Twin. The proposed solution integrates a directed-graph model with a multi-agent orchestration layer powered by CrewAI. The major functionalities implemented include real-time friction detection, in-memory lead time calculation, and autonomous agent-driven rerouting. The system achieved a significant response speedup of 11.95 hours compared to the human baseline, effectively meeting all project objectives and proving highly applicable for enhancing SME logistics resilience.

\section[Learning Outcomes]{Learning Outcomes}
The development and evaluation of this project have yielded the following key learning outcomes:

\begin{itemize}
    \item \textbf{Agentic AI Design:} Developed an understanding of how to architect multi-agent systems using CrewAI, including task chaining and context propagation to prevent LLM hallucinations.
    \item \textbf{Graph-Theoretic \& Inventory Modelling:} Gained proficiency in representing supply networks as directed graphs in NetworkX, applying Dijkstra's pathfinder and Chopra's Safety Stock framework to model volatility.
    \item \textbf{Systems Integration \& Software Practices:} Acquired experience in building client-server applications with FastAPI and React, while reinforcing modular design, separation of concerns, and fallback mechanisms.
    \item \textbf{Supply Chain Resilience Analytics:} Internalised the economic importance of maritime chokepoint resilience for Indian SMEs and the value of proactive, data-driven optimization tools.
\end{itemize}

\section[Limitations]{Limitations}
This section discusses the limitations and constraints of the proposed system:
\begin{itemize}
    \item \textbf{Strategic Scope:} The current implementation serves as a strategic digital twin focusing on high-level decision support rather than real-time tactical tracking.
    \item \textbf{Single-Shock Catering:} The system is currently optimised to process and mitigate one primary shock event per simulation cycle.
    \item \textbf{Regional Constraints:} The topology and business rules are specifically tailored for the Indian SME context, limiting immediate generalisability to other markets.
\end{itemize}

\section[Future Enhancement]{Future Enhancement}
Proposed improvements and future extensions for the system include:
\begin{itemize}
    \item \textbf{Dataset Expansion:} Increasing the richness of the dataset to include comprehensive multi-modal transport lanes.
    \item \textbf{Global Focus:} Expanding the target node focus beyond Indian SMEs to support international logistics corridors.
    \item \textbf{Operational Scaling:} Transitioning the framework into a fully operational digital twin by integrating live IoT telemetry and ERP data streams.
\end{itemize}
